% ---------------------------------------------------------------------------
% Packages and formatting. Kept to widely available packages so the document
% compiles identically on MiKTeX locally and on Overleaf.
% ---------------------------------------------------------------------------

\usepackage[T1]{fontenc}
\usepackage[utf8]{inputenc}
\usepackage{lmodern}

\usepackage[a4paper,left=1.5in,right=1in,top=1in,bottom=1in]{geometry}
\usepackage{setspace}
\onehalfspacing

% Lets TeX stretch a line a little further before giving up and letting text
% run into the margin. Mostly this absorbs the long file paths inside \todo
% placeholders while the report is still a draft.
\setlength{\emergencystretch}{3em}

\usepackage{graphicx}
\graphicspath{{figures/}}

\usepackage{amsmath}
\usepackage{amssymb}
\usepackage{booktabs}
\usepackage{longtable}
\usepackage{array}
\usepackage{enumitem}
\usepackage{xcolor}
\usepackage{caption}
\captionsetup{font=small,labelfont=bf}

\usepackage{fancyhdr}
\setlength{\headheight}{14pt}
\pagestyle{fancy}
\fancyhf{}
\fancyhead[L]{\small\projecttitle}
\fancyhead[R]{\small\thepage}
\renewcommand{\headrulewidth}{0.4pt}

% Chapter pages use the same header as the rest of the document.
\fancypagestyle{plain}{%
  \fancyhf{}%
  \fancyhead[L]{\small\projecttitle}%
  \fancyhead[R]{\small\thepage}%
  \renewcommand{\headrulewidth}{0.4pt}%
}

% Numeric, IEEE-like citations. unsrtnat ships with natbib, so there is
% nothing extra to install on either machine.
\usepackage[numbers,sort&compress]{natbib}
\bibliographystyle{unsrtnat}

\usepackage{listings}
\lstset{
  basicstyle=\ttfamily\footnotesize,
  breaklines=true,
  frame=single,
  captionpos=b,
  showstringspaces=false,
}

% Long reference URLs would otherwise run past the right margin, because a URL
% has no normal break points. xurl lets them break anywhere.
\usepackage{xurl}

% hyperref loads late by convention.
\usepackage[hidelinks]{hyperref}

% Shorthand for content that is still a placeholder. Every one of these
% should be gone before submission; \listtodos prints what is left.
\newcommand{\todo}[1]{\textcolor{red}{\textbf{[TODO: #1]}}}
